\section{Proofs}\label{apx:proofs}

\subsection*{Proof of Proposition \ref{prop:shapley_integral_identity_product_game}}
% \begin{proof}
By definition, the Shapley value of player $i\in\cD$ is
\[
\phi_i(v)
=
\sum_{S\subseteq \cD\setminus\{i\}}
\frac{|S|!\,(d-|S|-1)!}{d!}
\bigl(v(S\cup\{i\})-v(S)\bigr).
\]
For a product game, the marginal contribution simplifies to
\[
v(S\cup\{i\})-v(S)
=
\Bigl(\prod_{j\in S}u_j\Bigr)(u_i-1).
\]
Substituting this expression yields
\[
\phi_i(v)
=
(u_i-1)
\sum_{S\subseteq \cD\setminus\{i\}}
\mu_{|S|}
\prod_{j\in S}u_j,
\qquad
\mu_q := \frac{q!\,(d-q-1)!}{d!}.
\]

We now express the coefficients $\mu_q$ as integrals. For integers
$q\in\{0,\dots,d-1\}$, consider the integral
\[
\int_0^1 t^q(1-t)^{d-q-1}\,dt.
\]
By repeated integration by parts (or by the definition of the Beta function),
this integral evaluates to
\[
\int_0^1 t^q(1-t)^{d-q-1}\,dt
=
\frac{q!\,(d-q-1)!}{d!}
=
\mu_q.
\]
Thus, each Shapley weight $\mu_q$ admits the representation
\[
\mu_q
=
\int_0^1 t^q(1-t)^{d-q-1}\,dt.
\]

Substituting this identity into the Shapley expression and interchanging
summation and integration gives
\[
\phi_i(v)
=
(u_i-1)
\int_0^1
\sum_{S\subseteq \cD\setminus\{i\}}
t^{|S|}(1-t)^{d-|S|-1}
\prod_{j\in S}u_j
\,dt.
\]

To simplify the inner sum, observe that for each feature $j\neq i$, the term
inside the summation contributes either a factor $(1-t)$ if $j\notin S$ or
a factor $t\,u_j$ if $j\in S$. Summing over all subsets
$S\subseteq \cD\setminus\{i\}$ therefore corresponds to expanding the product
\[
\prod_{j\neq i}\bigl((1-t)+t\,u_j\bigr).
\]
Since $|\cD\setminus\{i\}|=d-1$, this expansion yields
\[
\sum_{S\subseteq \cD\setminus\{i\}}
t^{|S|}(1-t)^{d-|S|-1}\prod_{j\in S}u_j
=
\prod_{j\neq i}\bigl((1-t)+t\,u_j\bigr).
\]

Substituting back into the integral completes the proof:
\[
\phi_i(v)
=
(u_i-1)
\int_0^1
\prod_{j\neq i}\bigl((1-t)+t\,u_j\bigr)\,dt.
\]
% \end{proof}

\subsection*{Proof of Lemma~\ref{lem:gl_shapley_integral}}

\paragraph{Part~(i): Exactness.}
The integrand $g_i(t):=\prod_{j\neq i}\bigl((1-t)+tu_j\bigr)$ is a polynomial
of degree at most $d-1$ in $t$.
An $m_q$-point GL rule on $[0,1]$ integrates exactly every polynomial of degree
at most $2m_q-1$~\citep{davis1984methods}.
For $m_q\ge\lceil d/2\rceil$ we have $2m_q-1\ge 2\lceil d/2\rceil-1\ge d-1$,
so the quadrature sum reproduces the exact integral:
\[
\sum_{q=1}^{m_q}\omega_q\,g_i(\tau_q)
\;=\;
\int_0^1 g_i(t)\,dt.
\]
Multiplying both sides by $(u_i-1)$ and applying
Proposition~\ref{prop:shapley_integral_identity_product_game} gives
$\widehat{\phi}_i(v)=\phi_i(v)$.

\paragraph{Part~(ii): Approximation.}
Let
\[
g_i(t):=\prod_{j\neq i}\bigl((1-t)+t\,u_j\bigr).
\]
By Proposition~\ref{prop:shapley_integral_identity_product_game}, the exact
Shapley value admits the integral representation
\[
\phi_i(v)
=
(u_i-1)\int_0^1 g_i(t)\,dt,
\]
while its \(m_q\)-point Gauss--Legendre approximation is
\[
\widehat{\phi}_i(v)
=
(u_i-1)\sum_{q=1}^{m_q}\omega_q\,g_i(\tau_q).
\]
Therefore,
\[
\bigl|\phi_i(v)-\widehat{\phi}_i(v)\bigr|
=
|u_i-1|
\left|
\int_0^1 g_i(t)\,dt
-
\sum_{q=1}^{m_q}\omega_q\,g_i(\tau_q)
\right|.
\]

Now \(g_i\) is a polynomial in \(t\) of degree at most \(d-1\), hence it is
entire and in particular analytic on and inside every Bernstein ellipse
\(E_\rho\), for any \(\rho>1\). Standard Gauss--Legendre quadrature error
bounds for analytic functions on \([-1,1]\) \citep[Thm.~19.3]{trefethen2013atap},
transferred to \([0,1]\) via the affine change of variables \(x=2t-1\), imply
\[
\left|
\int_0^1 g_i(t)\,dt
-
\sum_{q=1}^{m_q}\omega_q\,g_i(\tau_q)
\right|
\le
\frac{4\max_{z\in E_\rho}|g_i(z)|}{\rho^{2m_q}(\rho^2-1)}.
\]
Substituting this bound into the previous display yields
\[
\bigl|\phi_i(v)-\widehat{\phi}_i(v)\bigr|
\le
|u_i-1|
\frac{4\max_{z\in E_\rho}|g_i(z)|}{\rho^{2m_q}(\rho^2-1)}.
\]
Hence, for any fixed \(\rho>1\),
\[
\bigl|\phi_i(v)-\widehat{\phi}_i(v)\bigr|
\le
C_\rho\,M_\rho\,|u_i-1|\,\rho^{-2m_q},
\]
where $M_\rho:=\max_{z\in E_\rho}|g_i(z)|$ and $C_\rho:=4/(\rho^2-1)$ depends only on $\rho$.
This proves the claimed geometric decay of the approximation error. \qed

\subsection*{Proof of Proposition~\ref{prop:complexity}}

\paragraph{Work.}
Computing, for each $q\in\{1,\dots,m_q\}$, the factors $T_{q,j}=(1-\tau_q)+\tau_q u_j$ and the associated quantities $\log|T_{q,j}|$ and $\mathrm{sgn}(T_{q,j})$ costs $O(d\,m_q)$ arithmetic operations in total. Accumulating $\log|P_q|=\sum_j \log|T_{q,j}|$ and $\mathrm{sgn}(P_q)=\prod_j \mathrm{sgn}(T_{q,j})$ is a length-$d$ reduction for each of the $m_q$ nodes, contributing $O(d\,m_q)$ additional work. Evaluating \cref{eq:gl_logspace_phi} then requires, for each $i\in\{1,\dots,d\}$, a weighted sum of $m_q$ terms, each obtained from $\log|P_q|-\log|T_{q,i}|$ in constant time, for a further $O(d\,m_q)$ operations. The total is therefore $O(d\,m_q)$.

\paragraph{Parallel time.}
Given an associative operator $\oplus$ and a length-$L$ sequence, a hierarchical parallel scan evaluates $a_1\oplus\cdots\oplus a_L$ in parallel time $O(\log L)$ using $O(L)$ processors~\citep{blelloch1990prefix}. Both $\log|P_q|$ (under addition) and $\mathrm{sgn}(P_q)$ (under multiplication) are such reductions over the feature axis of length $d$, and the $m_q$ nodes are independent, so the shared-product phase completes in parallel time $O(\log d)$. For each $i$, evaluating \cref{eq:gl_logspace_phi} requires (i) subtracting $\log|T_{q,i}|$ from $\log|P_q|$ pointwise (parallel time $O(1)$), and (ii) a weighted sum over $q$, which is a length-$m_q$ tree reduction with parallel time $O(\log m_q)$. Since the $d$ Shapley values are independent given $\{(\log|P_q|,\mathrm{sgn}(P_q))\}_{q=1}^{m_q}$, they are evaluated in parallel across $i$ without increasing parallel time. The overall parallel time is therefore $O(\log d+\log m_q)$. The required number of processors is the ratio of total work to parallel time~\citep{brent1974parallel}, giving $O\!\bigl(d\,m_q/(\log d+\log m_q)\bigr)$ processors. The exact regime $m_q=\lceil d/2\rceil$ is a special case.\qed

